\section{Metodologia dos Experimentos}

A automação dos experimentos é realizada pelo script \texttt{run\_experiment.py}, que executa cada configuração e coleta dados de tempo e telemetria do sistema de forma reprodutível.

\subsection{Parâmetros dos Experimentos}

A Tabela~\ref{tab:param_experimento} sumariza o protocolo adotado para cada caso de teste. O \textit{speedup} é calculado diretamente pela razão entre a média sequencial e a média paralela ($S_p = \overline{T}_{seq} / \overline{T}_p$).

\begin{table}[H]
\centering
\caption{Parâmetros e configurações de execução dos experimentos.}
\label{tab:param_experimento}
\begin{tabular}{@{}ll@{}}
\toprule
Parâmetro & Especificação / Valores \\ \midrule
Repetições por configuração & \textbf{20 vezes} consecutivas \\
Variantes testadas & Sequencial, OpenMP, MPI e Híbrido (MPI + OpenMP) \\
\textit{Threads} OpenMP & 1, 2, 4, 8 e 16 \\
Processos MPI & 1, 2, 4, 8 e 16 \\
Combinações Híbridas & $2\times2$, $2\times4$, $4\times2$, $2\times8$, $4\times4$ e $8\times2$ ($proc \times th$) \\
Tamanho do bloco (\texttt{BLOCK\_SIZE}) & \textbf{512} (OpenMP e Híbrido) \\
Tamanho do \textit{chunk} (\texttt{CHUNK\_SIZE}) & \textbf{1024} (MPI e Híbrido) \\
Afinidade de CPU & \texttt{OMP\_PROC\_BIND=close}, \texttt{OMP\_PLACES=cores} \\
Temporizador & \texttt{omp\_get\_wtime()} de alta precisão \\
Validação de corretude & Comparação estrita de \textit{score} contra a referência sequencial \\ \bottomrule
\end{tabular}
\end{table}

\subsection{Definição do \texttt{BLOCK\_SIZE}}

Para justificar a escolha do tamanho do bloco nas tarefas OpenMP, executamos uma varredura de sensibilidade com os valores 64, 128, 256 e 512 através do script \texttt{run\_test.py}. A Tabela~\ref{tab:blocksize} exibe o \textit{speedup} alcançado.

\begin{table}[H]
\centering
\caption{\textit{Speedup} da versão OpenMP em função do \texttt{BLOCK\_SIZE}.}
\label{tab:blocksize}
\begin{tabular}{@{}llccccc@{}}
\toprule
Entrada & Modo & \textit{Threads} & 64 & 128 & 256 & 512 \\ \midrule
\multirow{6}{*}{Q4\_random (u16)} & \multirow{3}{*}{\texttt{full-matrix}} & 4 & 1{,}55 & 1{,}77 & 1{,}90 & \textbf{2{,}13} \\
 & & 8 & 1{,}83 & 2{,}45 & 2{,}49 & \textbf{2{,}74} \\
 & & 16 & 3{,}10 & 2{,}84 & 3{,}05 & \textbf{3{,}21} \\ \cmidrule(l){2-7}
 & \multirow{3}{*}{\texttt{mem-opt}} & 4 & 1{,}93 & 2{,}85 & 3{,}43 & \textbf{3{,}62} \\
 & & 8 & 2{,}47 & 4{,}58 & 6{,}02 & \textbf{6{,}50} \\
 & & 16 & 3{,}07 & 9{,}26 & 10{,}73 & \textbf{11{,}39} \\ \midrule
\multirow{2}{*}{Q5\_random (u32)} & \multirow{2}{*}{\texttt{mem-opt}} & 8 & 2{,}02 & 3{,}09 & 3{,}76 & \textbf{4{,}00} \\
 & & 16 & 2{,}87 & 5{,}11 & 7{,}05 & \textbf{7{,}40} \\ \bottomrule
\end{tabular}
\end{table}

O tamanho de bloco 512 superou os demais em todas as configurações. Blocos pequenos (como 64) geram centenas de milhares de tarefas, sobrecarregando o escalonador do OpenMP. O valor 512 amortiza a sobrecarga de despacho mantendo granularidade suficiente para preencher os 16 núcleos.
